\documentclass[14pt,a4paper]{extarticle}

% ── Кодировка и язык ──────────────────────────────────────────────────────
\usepackage[utf8]{inputenc}
\usepackage[T2A]{fontenc}
\usepackage[russian]{babel}

% ── Поля страницы ─────────────────────────────────────────────────────────
\usepackage[left=30mm, right=15mm, top=20mm, bottom=20mm]{geometry}

% ── Математика ────────────────────────────────────────────────────────────
\usepackage{amsmath, amssymb}

% ── Таблицы ───────────────────────────────────────────────────────────────
\usepackage{booktabs}
\usepackage{array}
\usepackage{longtable}

% ── Графика ───────────────────────────────────────────────────────────────
\usepackage{graphicx}
\usepackage{float}

% ── Гиперссылки ───────────────────────────────────────────────────────────
\usepackage[hidelinks]{hyperref}

% ── Межстрочный интервал ──────────────────────────────────────────────────
\usepackage{setspace}
\onehalfspacing

% ── Отступы абзацев ───────────────────────────────────────────────────────
\usepackage{indentfirst}
\setlength{\parindent}{1.25cm}

% ── Нумерация разделов ────────────────────────────────────────────────────
\setcounter{secnumdepth}{2}

% ─────────────────────────────────────────────────────────────────────────
\begin{document}

% ── Титульный лист ────────────────────────────────────────────────────────
\begin{titlepage}
\begin{center}
\textbf{МИНИСТЕРСТВО ЦИФРОВОГО РАЗВИТИЯ, СВЯЗИ И МАССОВЫХ КОММУНИКАЦИЙ РОССИЙСКОЙ ФЕДЕРАЦИИ}

\vspace{0.3cm}
Ордена Трудового Красного Знамени федеральное государственное бюджетное учреждение высшего образования

<<Московский технический университет связи и информатики>>

\vspace{0.5cm}
Кафедра <<Математическая кибернетика и информационные технологии>>

\vspace{2cm}
{\Large \textbf{Курсовая работа}}

\vspace{0.3cm}
{\large <<Имитационное моделирование вычислительной системы>>}

\vspace{0.3cm}
{\large Вариант 22}

\vspace{0.3cm}
по дисциплине

<<Высоконагруженные приложения>>

\vspace{2cm}
\begin{flushright}
Выполнил:\\
студент группы БВТ2201 Шамсутдинов Р.\,Ф.\\[0.3cm]
Проверил:\\
преподаватель Александров А.\,Е.
\end{flushright}

\vfill
Москва, 2026 год
\end{center}
\end{titlepage}

% ── Аннотация ─────────────────────────────────────────────────────────────
\newpage
\section*{Аннотация}
\addcontentsline{toc}{section}{Аннотация}

В данной курсовой работе рассматривается задача имитационного моделирования
вычислительной системы, состоящей из трёх серверов с буфером на три программы.
Целью работы является исследование характеристик системы при различных законах
распределения входного потока и времени обслуживания заявок.

В работе разработана имитационная модель вычислительной системы с использованием
языка программирования Python. Модель реализована на основе метода
дискретно-событийного моделирования, что позволяет учитывать случайный характер
поступления заявок и времени их обработки. Проведено моделирование для двух типов
распределений: линейного и экспоненциального.

Для проверки корректности разработанной модели выполнена её верификация путём
сопоставления результатов имитационного моделирования с теоретическими значениями,
полученными на основе теории массового обслуживания. В качестве теоретической
модели рассмотрена система типа M/M/3/6 с ожиданием и ограниченным буфером, для
которой использованы формулы Эрланга.

В результате проведённых экспериментов получены все основные характеристики
системы: вероятности состояний $P_0$--$P_6$, относительная и абсолютная
пропускная способность, вероятность отказа, среднее число занятых серверов,
среднее число программ в системе и в буфере, среднее время нахождения программы
в системе и в буфере.

Установлено, что при экспоненциальном законе распределения результаты
имитационного моделирования хорошо согласуются с теоретическими значениями.
При линейном распределении система работает в условиях значительно более высокой
нагрузки ($\rho = 4{,}5$ против $\rho = 3{,}0$), что приводит к большей
вероятности отказа и более длительному времени ожидания.

% ── Содержание ────────────────────────────────────────────────────────────
\newpage
\tableofcontents

% ─────────────────────────────────────────────────────────────────────────
\newpage
\section{Введение}

\subsection{Актуальность темы}

Современные вычислительные системы, используемые в высоконагруженных
приложениях, обрабатывают значительные объёмы запросов, поступающих в случайные
моменты времени. Эффективность работы таких систем напрямую зависит от их
способности обрабатывать входной поток заявок без существенных задержек и потерь.

Одним из основных инструментов исследования подобных систем является теория
массового обслуживания, позволяющая анализировать характеристики систем с учётом
случайного характера входных потоков и времени обслуживания. Однако аналитические
методы применимы лишь к ограниченному числу моделей и зачастую требуют
упрощающих предположений, таких как экспоненциальное распределение. В связи с
этим особую актуальность приобретает имитационное моделирование, позволяющее
исследовать поведение вычислительных систем при произвольных законах
распределения и сложных условиях функционирования.

\subsection{Цель и задачи работы}

Целью данной курсовой работы является разработка имитационной модели
вычислительной системы с буфером и исследование её характеристик при различных
законах распределения входного потока и времени обслуживания заявок.

Для достижения поставленной цели необходимо решить следующие задачи:
\begin{itemize}
  \item разработать имитационную модель вычислительной системы с тремя серверами
        и буфером на три программы;
  \item реализовать модель с использованием метода дискретно-событийного
        моделирования;
  \item провести моделирование системы при линейном и экспоненциальном законах
        распределения;
  \item определить все характеристики системы: вероятности состояний $P_0$--$P_6$,
        пропускную способность, вероятность отказа, среднее число занятых серверов,
        среднее число программ в системе и в буфере, среднее время нахождения
        программы в системе и в буфере;
  \item разработать теоретическую модель на основе теории массового обслуживания
        (M/M/3/6);
  \item выполнить верификацию имитационной модели путём сопоставления с
        теоретическими результатами;
  \item провести анализ результатов машинных экспериментов и сделать выводы о
        влиянии закона распределения на характеристики системы.
\end{itemize}

\subsection{Выбор языка программирования}

Для реализации имитационной модели выбран язык программирования Python. Данный
выбор обусловлен рядом преимуществ.

Во-первых, Python обладает простым и понятным синтаксисом, что позволяет быстро
разрабатывать и модифицировать модели. Во-вторых, язык предоставляет широкий
набор библиотек для научных вычислений и обработки данных: \texttt{NumPy} ---
для генерации случайных величин, \texttt{Pandas} --- для обработки результатов,
\texttt{Matplotlib} --- для построения графиков.

Кроме того, Python хорошо подходит для реализации дискретно-событийных моделей
благодаря наличию встроенных структур данных, таких как приоритетные очереди
(\texttt{heapq}), что позволяет эффективно организовать обработку событий во
времени.

\subsection{Объект исследования}

Объектом исследования в данной работе является вычислительная система,
состоящая из трёх параллельно работающих серверов, предназначенных для обработки
поступающих программ (заявок), и буфера ёмкостью три программы.

При поступлении программы она направляется на первый свободный сервер. Если все
серверы заняты, программа помещается в буфер (очередь ожидания). Если буфер
также заполнен, программа покидает систему необработанной. Это соответствует
системе массового обслуживания с ожиданием и ограниченной ёмкостью.

В рамках исследования рассматриваются два варианта входного потока и времени
обслуживания: линейное и экспоненциальное распределения.

\subsection{Структура работы}

Курсовая работа состоит из введения, пяти основных разделов, заключения, списка
использованных источников и приложения.

Во введении обосновывается актуальность темы, формулируются цель и задачи работы.
Во втором разделе рассматривается имитационная модель вычислительной системы,
включая вероятностные характеристики входного потока, логическую схему модели и
принципы управления модельным временем. Третий раздел посвящён описанию
теоретической модели системы на основе теории массового обслуживания. В четвёртом
разделе проводится верификация имитационной модели. Пятый раздел содержит анализ
результатов машинных экспериментов. В заключении приводятся основные выводы.

% ─────────────────────────────────────────────────────────────────────────
\newpage
\section{Имитационная модель вычислительной системы}

\subsection{Вероятностные характеристики входного потока и обработки заявок}

В рассматриваемой вычислительной системе входной поток заявок и время их
обслуживания являются случайными величинами. Для их описания используются два
различных закона распределения: линейный и экспоненциальный.

\subsubsection*{Линейное распределение}

В первом случае интервалы между поступлениями заявок распределены по линейному
(равномерному) закону на отрезке $[T_{z\min},\, T_{z\max}]$, где:
\[
  T_{z\min} = \tfrac{1}{2}\ \text{с}, \quad T_{z\max} = \tfrac{5}{6}\ \text{с}.
\]

Плотность распределения:
\[
  f(t) = \frac{1}{T_{z\max} - T_{z\min}} = \frac{1}{\,5/6 - 1/2\,} = 3,
  \quad t \in \bigl[\tfrac{1}{2};\, \tfrac{5}{6}\bigr].
\]

Математическое ожидание (средний интервал между приходами):
\[
  \mathbb{E}[T_z] = \frac{T_{z\min} + T_{z\max}}{2}
  = \frac{1/2 + 5/6}{2} = \frac{2}{3} \approx 0{,}667\ \text{с}.
\]

Интенсивность потока:
\[
  \lambda = \frac{1}{\mathbb{E}[T_z]} = \frac{1}{2/3} = 1{,}5\ \text{[1/с]}.
\]

Время обслуживания заявок распределено равномерно на интервале
$[T_{s\min},\, T_{s\max}]$:
\[
  T_{s\min} = 1\ \text{с}, \quad T_{s\max} = 5\ \text{с}.
\]

Математическое ожидание времени обслуживания:
\[
  \mathbb{E}[T_s] = \frac{T_{s\min} + T_{s\max}}{2} = \frac{1 + 5}{2} = 3\ \text{с}.
\]

Интенсивность обслуживания одним сервером:
\[
  \mu = \frac{1}{\mathbb{E}[T_s]} = \frac{1}{3} \approx 0{,}333\ \text{[1/с]}.
\]

Нагрузка на систему:
\[
  \rho = \frac{\lambda}{\mu} = \frac{1{,}5}{1/3} = 4{,}5.
\]

Линейное распределение характеризуется ограниченной вариативностью, что означает
отсутствие резких скачков в интенсивности входного потока.

\subsubsection*{Экспоненциальное распределение}

Во втором случае интервалы между поступлениями заявок подчиняются
экспоненциальному закону с параметром $\lambda = 1{,}5$~[1/с]:
\[
  f(t) = \lambda e^{-\lambda t} = 1{,}5\, e^{-1{,}5\, t}, \quad t \geq 0.
\]

Математическое ожидание:
\[
  \mathbb{E}[T_z] = \frac{1}{\lambda} = \frac{1}{1{,}5} \approx 0{,}667\ \text{с}.
\]

Время обслуживания также распределено экспоненциально. По условию среднее время
обслуживания $t_{\text{обр}} = 2$~с, следовательно:
\[
  \mu = \frac{1}{t_{\text{обр}}} = \frac{1}{2} = 0{,}5\ \text{[1/с]}.
\]

Нагрузка на систему:
\[
  \rho = \frac{\lambda}{\mu} = \frac{1{,}5}{0{,}5} = 3{,}0.
\]

Экспоненциальное распределение обладает свойством отсутствия памяти, что делает
его ключевым при построении моделей теории массового обслуживания.

\subsubsection*{Сравнение распределений}

Основное различие между линейным и экспоненциальным распределениями заключается
в среднем времени обслуживания и, как следствие, в нагрузке на систему:
\begin{itemize}
  \item линейное распределение: $\mathbb{E}[T_s] = 3$~с, $\rho = 4{,}5$ ---
        высокая нагрузка, система работает в режиме перегрузки;
  \item экспоненциальное распределение: $\mathbb{E}[T_s] = 2$~с, $\rho = 3{,}0$ ---
        умеренная нагрузка, система работает в сбалансированном режиме.
\end{itemize}

\subsection{Логическая схема модели вычислительной системы и её описание}

Рассматриваемая вычислительная система состоит из трёх параллельно работающих
серверов и буфера ёмкостью три программы. Максимальное число программ в системе
одновременно равно шести.

\subsubsection*{Принцип работы системы}

При поступлении программы выполняется проверка состояния серверов:
\begin{itemize}
  \item если имеется хотя бы один свободный сервер, программа направляется на него;
  \item если все серверы заняты, но буфер не заполнен, программа помещается в
        буфер (очередь FIFO);
  \item если все серверы заняты и буфер заполнен, программа получает отказ и
        покидает систему;
  \item после завершения обработки сервер освобождается; если в буфере есть
        программы, первая из них немедленно начинает обслуживаться.
\end{itemize}

\subsubsection*{Состояния системы}

Система имеет 7 состояний, представленных в таблице~\ref{tab:states}.

\begin{table}[H]
\centering
\caption{Состояния вычислительной системы}
\label{tab:states}
\begin{tabular}{clc}
\toprule
Состояние & Описание & Программ в буфере \\
\midrule
$P_0$ & ВС не загружена & 0 \\
$P_1$ & Загружен 1 сервер & 0 \\
$P_2$ & Загружены 2 сервера & 0 \\
$P_3$ & Загружены 3 сервера, буфер пуст & 0 \\
$P_4$ & Загружены 3 сервера & 1 \\
$P_5$ & Загружены 3 сервера & 2 \\
$P_6$ & Загружены 3 сервера, буфер заполнен & 3 \\
\bottomrule
\end{tabular}
\end{table}

\subsubsection*{Описание алгоритма моделирования}

Имитационная модель реализована на основе дискретно-событийного подхода.
Основными типами событий являются:
\begin{itemize}
  \item \textbf{arrival} --- поступление новой программы;
  \item \textbf{departure} --- завершение обслуживания программы.
\end{itemize}

Для хранения событий используется приоритетная очередь (\texttt{heapq}),
упорядоченная по времени наступления событий. Это позволяет в каждый момент
времени обрабатывать ближайшее событие.

Алгоритм работы модели включает следующие этапы:

\begin{enumerate}
  \item \textbf{Инициализация системы:} задаются параметры модели (параметры
        распределений, время моделирования, число серверов, ёмкость буфера);
        генерируется момент первого поступления программы; очередь событий
        инициализируется первым событием.

  \item \textbf{Обработка событий:} выбирается ближайшее событие из очереди;
        вычисляется интервал времени $\Delta t$ с момента предыдущего события;
        обновляется статистика состояний системы; в зависимости от типа события
        выполняются соответствующие действия.

  \item \textbf{Обработка поступления программы:}
        \begin{itemize}
          \item увеличивается общее число программ;
          \item генерируется следующее поступление;
          \item если есть свободный сервер --- программа принимается, планируется
                событие завершения обслуживания;
          \item если серверы заняты, но буфер не полон --- программа помещается
                в буфер, фиксируется время прихода в буфер;
          \item иначе --- фиксируется отказ.
        \end{itemize}

  \item \textbf{Обработка завершения обслуживания:}
        \begin{itemize}
          \item если в буфере есть программы --- первая из них немедленно
                начинает обслуживаться; фиксируется время ожидания в буфере;
          \item иначе --- сервер освобождается.
        \end{itemize}

  \item \textbf{Завершение моделирования:} моделирование продолжается до
        достижения заданного времени $T = 3600$~с; вычисляются итоговые
        характеристики системы.
\end{enumerate}

\subsubsection*{Сбор статистики}

В процессе моделирования рассчитываются следующие показатели:
\begin{itemize}
  \item вероятности состояний системы $P_0, P_1, \ldots, P_6$;
  \item относительная пропускная способность:
        $Q = N_{\text{обсл}} / N_{\text{вход}}$;
  \item абсолютная пропускная способность:
        $A = N_{\text{обсл}} / T$;
  \item вероятность отказа:
        $P_{\text{отк}} = N_{\text{отказ}} / N_{\text{вход}}$;
  \item среднее число занятых серверов:
        $K = \frac{1}{T}\int_0^T k(t)\,dt$;
  \item среднее число программ в буфере:
        $N_{\text{буф}} = \frac{1}{T}\int_0^T b(t)\,dt$;
  \item среднее число программ в системе:
        $N_{\text{прог}} = \sum_{i=0}^{6} i \cdot P_i$;
  \item среднее время нахождения программы в системе (формула Литтла):
        $T_{\text{прог}} = N_{\text{прог}} / A$;
  \item среднее время ожидания в буфере:
        $T_{\text{буф}} = \sum w_i / N_{\text{буф,прошло}}$,
        где $w_i$ --- время ожидания $i$-й программы в буфере.
\end{itemize}

\subsection{Управление модельным временем}

В разработанной модели используется метод дискретно-событийного моделирования,
при котором время изменяется не равномерно, а скачкообразно --- от одного события
к другому:
\[
  t_{\text{текущее}} = t_{\text{следующего события}}.
\]

Это означает, что модель <<перескакивает>> через промежутки времени, в которых не
происходит никаких изменений состояния системы.

\subsubsection*{Преимущества метода}

Использование моделирования по особым событиям обладает рядом преимуществ:
\begin{itemize}
  \item высокая вычислительная эффективность, так как обрабатываются только
        значимые моменты времени;
  \item отсутствие накопления ошибок, связанных с дискретизацией времени;
  \item точное соответствие реальным процессам обслуживания.
\end{itemize}

\subsubsection*{Учёт времени в состояниях}

Для вычисления вероятностей состояний системы используется накопление времени
пребывания в каждом состоянии. Если в течение интервала $\Delta t$ система
находилась в состоянии $k$, то:
\[
  T_k \mathrel{+}= \Delta t.
\]

После завершения моделирования вероятности состояний определяются как:
\[
  P_k = \frac{T_k}{T}.
\]

Среднее число занятых серверов и среднее число программ в буфере вычисляются
методом численного интегрирования по времени:
\[
  K = \frac{1}{T}\sum_i k_i \cdot \Delta t_i, \qquad
  N_{\text{буф}} = \frac{1}{T}\sum_i b_i \cdot \Delta t_i,
\]
где $k_i$ и $b_i$ --- число занятых серверов и число программ в буфере на
$i$-м интервале между событиями.

% ─────────────────────────────────────────────────────────────────────────
\newpage
\section{Модель ВС на основе теории массового обслуживания}

\subsection{Описание упрощённой структуры ВС на основе СМО}

Для анализа работы вычислительной системы наряду с имитационной моделью
используется её упрощённое описание на основе теории массового обслуживания.
Такой подход позволяет получить аналитические зависимости для основных
характеристик системы и использовать их для верификации результатов моделирования.

\subsubsection*{Тип системы массового обслуживания}

Рассматриваемая вычислительная система представляется в виде системы массового
обслуживания типа \textbf{M/M/3/6}, где:
\begin{itemize}
  \item первый символ M --- экспоненциальный закон распределения интервалов между
        поступлениями заявок;
  \item второй символ M --- экспоненциальное распределение времени обслуживания;
  \item число 3 --- количество каналов обслуживания (серверов);
  \item число 6 --- максимальное число заявок в системе (3 сервера + 3 в буфере).
\end{itemize}

\subsubsection*{Основные параметры системы}

\[
  \lambda = 1{,}5\ \text{[1/с]}, \quad
  \mu = 0{,}5\ \text{[1/с]}, \quad
  \rho = \frac{\lambda}{\mu} = \frac{1{,}5}{0{,}5} = 3{,}0.
\]

\subsubsection*{Вероятности состояний системы}

Вероятность того, что в системе находится $n$ программ, определяется по формулам
Эрланга для M/M/c/K:

\[
  P_n = \begin{cases}
    P_0 \cdot \dfrac{\rho^n}{n!}, & 0 \le n \le c, \\[8pt]
    P_0 \cdot \dfrac{\rho^n}{c!\, c^{n-c}}, & c < n \le K,
  \end{cases}
\]

где $c = 3$, $K = 6$, а $P_0$ находится из условия нормировки:

\[
  P_0 = \left(\sum_{n=0}^{c} \frac{\rho^n}{n!}
             + \sum_{n=c+1}^{K} \frac{\rho^n}{c!\, c^{n-c}}\right)^{-1}.
\]

\subsubsection*{Основные характеристики системы}

\begin{enumerate}
  \item \textbf{Вероятность отказа.} Отказ происходит при полной загрузке системы
        (состояние $n = K = 6$):
        \[
          P_{\text{отк}} = P_6.
        \]

  \item \textbf{Относительная пропускная способность:}
        \[
          Q = 1 - P_{\text{отк}}.
        \]

  \item \textbf{Абсолютная пропускная способность:}
        \[
          A = \lambda \cdot Q.
        \]

  \item \textbf{Среднее число занятых серверов:}
        \[
          K = \sum_{n=0}^{K} \min(n,\, c) \cdot P_n.
        \]

  \item \textbf{Среднее число программ в буфере:}
        \[
          N_{\text{буф}} = \sum_{n=c+1}^{K} (n - c) \cdot P_n.
        \]

  \item \textbf{Среднее число программ в системе:}
        \[
          N_{\text{прог}} = \sum_{n=0}^{K} n \cdot P_n.
        \]

  \item \textbf{Среднее время нахождения программы в системе и в буфере}
        (формула Литтла $N = A \cdot T$):
        \[
          T_{\text{прог}} = \frac{N_{\text{прог}}}{A}, \qquad
          T_{\text{буф}}  = \frac{N_{\text{буф}}}{A}.
        \]
\end{enumerate}

\subsubsection*{Особенности теоретической модели}

Теоретическая модель основывается на ряде допущений:
\begin{itemize}
  \item входной поток является простейшим (пуассоновским);
  \item время обслуживания имеет экспоненциальное распределение;
  \item процессы поступления и обслуживания независимы.
\end{itemize}

Эти предположения позволяют получить аналитическое решение, однако ограничивают
применимость модели. В частности, при использовании линейного закона
распределения аналитические методы становятся неприменимыми.

Таким образом, теоретическая модель служит эталоном для проверки корректности
имитационной модели в случае экспоненциальных распределений.

% ─────────────────────────────────────────────────────────────────────────
\newpage
\section{Верификация и проверка достоверности логической схемы модели}

\subsection{Сопоставление ИМ с упрощённой структурой ВС на основе СМО}

Одним из важнейших этапов разработки имитационной модели является её верификация,
то есть проверка соответствия полученных результатов известным аналитическим
решениям. В данной работе верификация проводится путём сопоставления результатов
имитационного моделирования с теоретическими значениями, полученными на основе
модели M/M/3/6.

\subsubsection*{Методика верификации}

Для проверки корректности модели выполняются следующие действия:
\begin{itemize}
  \item задаются параметры системы, соответствующие условиям теоретической модели:
        экспоненциальное распределение входного потока ($\lambda = 1{,}5$~[1/с])
        и экспоненциальное распределение времени обслуживания ($\mu = 0{,}5$~[1/с]);
  \item проводится серия из 10 имитационных экспериментов с последующим
        усреднением результатов;
  \item рассчитываются теоретические значения характеристик по формулам Эрланга;
  \item выполняется сравнение полученных результатов.
\end{itemize}

\subsubsection*{Сравнение результатов}

Результаты сравнения имитационной и теоретической моделей приведены в
таблице~\ref{tab:verif}.

\begin{table}[H]
\centering
\caption{Сравнение характеристик теоретической и имитационной модели (10 запусков)}
\label{tab:verif}
\begin{tabular}{lcc}
\toprule
Характеристика & Имитация (эксп) & Теория M/M/3/6 \\
\midrule
$P_0$ & 0{,}037742 & 0{,}037736 \\
$P_1$ & 0{,}114785 & 0{,}113208 \\
$P_2$ & 0{,}171803 & 0{,}169811 \\
$P_3$ & 0{,}171454 & 0{,}169811 \\
$P_4$ & 0{,}168907 & 0{,}169811 \\
$P_5$ & 0{,}168180 & 0{,}169811 \\
$P_6$ & 0{,}167129 & 0{,}169811 \\
$Q$   & 0{,}831808 & 0{,}830189 \\
$A$, [1/с] & 1{,}242556 & 1{,}245283 \\
$P_{\text{отк}}$ & 0{,}167430 & 0{,}169811 \\
$K$   & 2{,}485400 & 2{,}490566 \\
$N_{\text{прог}}$ & 3{,}492054 & 3{,}509434 \\
$T_{\text{прог}}$, с & 2{,}811200 & 2{,}818182 \\
$N_{\text{буф}}$ & 1{,}006653 & 1{,}018868 \\
$T_{\text{буф}}$, с & 1{,}328829 & 0{,}818182 \\
\bottomrule
\end{tabular}
\end{table}

Из представленных данных видно, что значения, полученные в имитационной модели
при экспоненциальном распределении, практически совпадают с теоретическими.
Различия между значениями не превышают 1--2\%, что свидетельствует о высокой
точности моделирования.

Небольшие отклонения обусловлены следующими факторами:
\begin{itemize}
  \item случайный характер генерации входного потока и времени обслуживания;
  \item конечное время моделирования (1 час);
  \item статистические погрешности, возникающие при имитации.
\end{itemize}

Небольшое расхождение $T_{\text{буф}}$ объясняется тем, что в имитационной
модели $T_{\text{буф}}$ измеряется прямым способом (среднее время ожидания
программ, прошедших через буфер), тогда как в теоретической модели используется
формула Литтла $T_{\text{буф}} = N_{\text{буф}} / A$.

\subsubsection*{Сравнение вероятностей состояний}

Графическое сравнение вероятностей состояний $P_0$--$P_6$ показывает
практически полное совпадение кривых для имитационной и теоретической моделей.
Это означает, что распределение состояний системы, полученное в ходе
моделирования, соответствует стационарному распределению, рассчитанному по
формулам Эрланга.

Совпадение подтверждает корректность:
\begin{itemize}
  \item логической схемы модели;
  \item алгоритма обработки событий;
  \item механизма учёта времени в состояниях.
\end{itemize}

\subsubsection*{Выводы по верификации}

Проведённая верификация позволяет сделать следующие выводы:
\begin{itemize}
  \item имитационная модель при экспоненциальном распределении корректно
        воспроизводит поведение системы массового обслуживания;
  \item разработанный алгоритм дискретно-событийного моделирования реализован
        правильно;
  \item полученные характеристики системы соответствуют теоретическим значениям.
\end{itemize}

Таким образом, экспоненциальный закон соответствует предположениям теории
массового обслуживания, а разработанная модель является достоверной и может
использоваться для дальнейшего анализа.

% ─────────────────────────────────────────────────────────────────────────
\newpage
\section{Анализ результатов машинных экспериментов}

На завершающем этапе работы проведён анализ результатов имитационного
моделирования вычислительной системы при двух различных законах распределения:
линейном и экспоненциальном. Целью анализа является сравнение характеристик
системы и определение влияния закона распределения на её поведение. Для этого
проводилась серия из 10 имитационных экспериментов для каждого вида распределения
с последующим усреднением результатов.

\subsection{Сравнение основных характеристик системы}

Результаты моделирования представлены в таблице~\ref{tab:compare}.

\begin{table}[H]
\centering
\caption{Сравнение характеристик ВС для линейного и экспоненциального распределений (10 запусков)}
\label{tab:compare}
\begin{tabular}{lcc}
\toprule
Характеристика & Линейный закон & Экспоненциальный закон \\
\midrule
$P_0$ & 0{,}000193 & 0{,}037742 \\
$P_1$ & 0{,}000187 & 0{,}114785 \\
$P_2$ & 0{,}000408 & 0{,}171803 \\
$P_3$ & 0{,}005837 & 0{,}171454 \\
$P_4$ & 0{,}089069 & 0{,}168907 \\
$P_5$ & 0{,}370024 & 0{,}168180 \\
$P_6$ & 0{,}534283 & 0{,}167129 \\
$Q$   & 0{,}666281 & 0{,}831808 \\
$A$, [1/с] & 0{,}999306 & 1{,}242556 \\
$P_{\text{отк}}$ & 0{,}332664 & 0{,}167430 \\
$K$   & 2{,}998640 & 2{,}485400 \\
$N_{\text{прог}}$ & 5{,}430606 & 3{,}492054 \\
$T_{\text{прог}}$, с & 5{,}434731 & 2{,}811200 \\
$N_{\text{буф}}$ & 2{,}431966 & 1{,}006653 \\
$T_{\text{буф}}$, с & 2{,}434817 & 1{,}328829 \\
\bottomrule
\end{tabular}
\end{table}

\subsection{Анализ вероятностей состояний}

На основе полученных данных видно, что распределение вероятностей состояний
системы существенно различается для двух законов распределения.

\textbf{При линейном распределении} ($\rho = 4{,}5$):
\begin{itemize}
  \item вероятность полной загрузки системы $P_6 \approx 0{,}534$ --- буфер
        заполнен более половины времени;
  \item вероятности состояний с малым числом занятых серверов
        ($P_0, P_1, P_2$) близки к нулю --- система практически никогда не
        бывает свободной;
  \item система большую часть времени находится в состоянии перегрузки.
\end{itemize}

\textbf{При экспоненциальном распределении} ($\rho = 3{,}0$):
\begin{itemize}
  \item вероятности состояний $P_3$--$P_6$ примерно одинаковы ($\approx 0{,}17$)
        --- это математически закономерный результат при $\rho = c = 3$;
  \item система чаще находится в состояниях частичной загрузки;
  \item наблюдается баланс между поступлением и обслуживанием заявок.
\end{itemize}

\subsection{Анализ основных характеристик}

\textbf{Относительная пропускная способность:}
\[
  Q_{\text{лин}} = 0{,}666, \quad Q_{\text{эксп}} = 0{,}832.
\]
Экспоненциальное распределение обеспечивает более высокую вероятность
обслуживания заявок. Это объясняется меньшим средним временем обслуживания
(2~с против 3~с), что снижает нагрузку на систему.

\textbf{Абсолютная пропускная способность:}
\[
  A_{\text{лин}} = 0{,}999\ \text{[1/с]}, \quad
  A_{\text{эксп}} = 1{,}243\ \text{[1/с]}.
\]
Абсолютная пропускная способность также выше при экспоненциальном распределении.
Система обрабатывает на 24\% больше программ в единицу времени.

\textbf{Вероятность отказа:}
\[
  P_{\text{отк,лин}} = 0{,}333, \quad P_{\text{отк,эксп}} = 0{,}167.
\]
При линейном распределении каждая третья программа получает отказ, при
экспоненциальном --- каждая шестая. Это связано с тем, что при линейном
распределении система работает в режиме постоянной высокой загрузки.

\textbf{Среднее число занятых серверов:}
\[
  K_{\text{лин}} = 2{,}999 \approx 3, \quad K_{\text{эксп}} = 2{,}485.
\]
При линейном распределении все три сервера заняты практически постоянно.
При экспоненциальном распределении нагрузка распределяется менее равномерно,
что приводит к снижению среднего числа занятых серверов.

\textbf{Среднее число программ в системе и время нахождения:}
\[
  N_{\text{прог,лин}} = 5{,}43, \quad N_{\text{прог,эксп}} = 3{,}49.
\]
\[
  T_{\text{прог,лин}} = 5{,}43\ \text{с}, \quad T_{\text{прог,эксп}} = 2{,}81\ \text{с}.
\]
Программа при линейном законе проводит в системе почти вдвое больше времени,
чем при экспоненциальном. Это объясняется большим временем ожидания в буфере.

\textbf{Среднее число программ в буфере и время ожидания:}
\[
  N_{\text{буф,лин}} = 2{,}43, \quad N_{\text{буф,эксп}} = 1{,}01.
\]
\[
  T_{\text{буф,лин}} = 2{,}43\ \text{с}, \quad T_{\text{буф,эксп}} = 1{,}33\ \text{с}.
\]
При линейном законе в буфере в среднем находится 2{,}4 программы, тогда как
при экспоненциальном --- около 1. Время ожидания в буфере при линейном законе
почти вдвое больше.

\subsection{Интерпретация результатов}

Полученные результаты демонстрируют принципиальное различие между двумя типами
распределений:
\begin{itemize}
  \item \textbf{линейное распределение} создаёт устойчивую и плотную нагрузку
        на систему ($\rho = 4{,}5$), приводящую к её перегрузке: буфер заполнен
        более половины времени, каждая третья программа получает отказ;
  \item \textbf{экспоненциальное распределение} формирует случайный поток с
        высокой вариативностью ($\rho = 3{,}0$), что позволяет системе
        периодически разгружаться.
\end{itemize}

Несмотря на то, что средние значения интервалов между заявками в обоих случаях
одинаковы ($\mathbb{E}[T_z] \approx 0{,}667$~с), различие в среднем времени
обслуживания (3~с против 2~с) существенно влияет на все характеристики системы.

В предыдущем разделе было показано, что при экспоненциальном распределении
результаты имитационной модели практически совпадают с теоретическими значениями
для системы M/M/3/6. В то же время линейное распределение не удовлетворяет
предположениям классической теории СМО, что приводит к значительным отличиям
в характеристиках системы.

% ─────────────────────────────────────────────────────────────────────────
\newpage
\section*{Заключение}
\addcontentsline{toc}{section}{Заключение}

В ходе выполнения курсовой работы была разработана имитационная модель
вычислительной системы, состоящей из трёх серверов с буфером на три программы.
Модель реализована с использованием метода дискретно-событийного моделирования
на языке программирования Python.

В рамках работы были исследованы два варианта входного потока и времени
обслуживания заявок: линейное и экспоненциальное распределения. Для оценки
корректности модели проведена её верификация путём сопоставления с теоретической
моделью системы массового обслуживания типа M/M/3/6.

Результаты показали, что при экспоненциальном распределении характеристики
имитационной модели практически совпадают с теоретическими значениями. Это
подтверждает корректность разработанной модели и её соответствие классическим
положениям теории массового обслуживания.

Анализ машинных экспериментов показал, что закон распределения оказывает
существенное влияние на поведение системы. При линейном распределении
($\rho = 4{,}5$) наблюдается более высокая загрузка серверов ($K \approx 3$),
большая вероятность отказа ($P_{\text{отк}} \approx 0{,}33$) и более длительное
время нахождения программы в системе ($T_{\text{прог}} \approx 5{,}4$~с).
При экспоненциальном распределении ($\rho = 3{,}0$) система работает более
эффективно: $P_{\text{отк}} \approx 0{,}17$, $T_{\text{прог}} \approx 2{,}8$~с.

Главный вывод работы заключается в том, что экспоненциальное распределение
адекватно описывает функционирование систем массового обслуживания и согласуется
с их теоретическими моделями. В то же время линейное распределение не
соответствует предположениям СМО и приводит к иным характеристикам системы.

Таким образом, использование имитационного моделирования позволяет не только
подтвердить теоретические результаты, но и исследовать поведение систем в
условиях, когда аналитические методы неприменимы.

% ─────────────────────────────────────────────────────────────────────────
\newpage
\section*{Список использованных источников}
\addcontentsline{toc}{section}{Список использованных источников}

\begin{enumerate}
  \item Боев, В.\,Д. Имитационное моделирование систем: учебник для
        академического бакалавриата / В.\,Д. Боев. --- М.: Юрайт, 2019.
  \item Карпов, Ю.\,Г. Имитационное моделирование систем. Введение в
        моделирование с AnyLogic и Python / Ю.\,Г. Карпов. --- СПб.:
        БХВ-Петербург, 2021.
  \item Рыжиков, Ю.\,И. Имитационное моделирование. Теория и технологии /
        Ю.\,И. Рыжиков. --- СПб.: Профессиональная литература, 2020.
  \item Кельтон, В. Имитационное моделирование. Классика CS / В. Кельтон,
        А. Лоу. --- 3-е изд. --- СПб.: Питер, 2024.
  \item Власов, С.\,Е. Моделирование вычислительных систем и сетей: учебное
        пособие / С.\,Е. Власов. --- М.: Инфра-М, 2022.
  \item Banks, J. Discrete-Event System Simulation / J. Banks, J.\,S. Carson,
        B.\,L. Nelson, D.\,M. Nicol. --- 6th ed. --- Pearson, 2019.
  \item Fishman, G.\,S. Discrete-Event Simulation: Modeling, Programming, and
        Analysis / G.\,S. Fishman. --- Springer, 2021.
  \item Zeigler, B.\,P. Theory of Modeling and Simulation / B.\,P. Zeigler,
        A. Muzy, E. Kofman. --- 3rd ed. --- Academic Press, 2018.
\end{enumerate}

% ─────────────────────────────────────────────────────────────────────────
\newpage
\section*{Приложение}
\addcontentsline{toc}{section}{Приложение}

\subsection*{Исходный код имитационной модели (Python)}

\begin{verbatim}
import math
import numpy as np
import pandas as pd
import heapq
import matplotlib.pyplot as plt
from collections import deque

np.random.seed(42)


def linear_time(t_min, t_max):
    return np.random.uniform(t_min, t_max)


def exponential_time(rate):
    return np.random.exponential(1.0 / rate)


def simulate_system(
    arrival_type="linear", service_type="linear",
    T=3600, n_servers=3, buffer_size=3,
    tz_min=1/2, tz_max=5/6, ts_min=1.0, ts_max=5.0,
    lmbda=1.5, mu=0.5
):
    max_in_system = n_servers + buffer_size
    n_states = max_in_system + 1

    current_time = 0.0
    last_event_time = 0.0
    busy_servers = 0
    buffer_count = 0
    events = []

    total_arrivals = 0
    total_served = 0
    total_rejected = 0

    state_time = np.zeros(n_states)
    busy_servers_sum = 0.0
    buffer_sum = 0.0
    total_wait_in_buffer = 0.0
    total_passed_through_buffer = 0
    buffer_queue = deque()

    def gen_arrival():
        if arrival_type == "linear":
            return linear_time(tz_min, tz_max)
        else:
            return exponential_time(lmbda)

    def gen_service():
        if service_type == "linear":
            return linear_time(ts_min, ts_max)
        else:
            return exponential_time(mu)

    heapq.heappush(events, (gen_arrival(), "arrival", None))

    while events:
        current_time, event_type, event_data = heapq.heappop(events)
        if current_time > T:
            break

        dt = current_time - last_event_time
        state = busy_servers + buffer_count
        state_time[state] += dt
        busy_servers_sum += busy_servers * dt
        buffer_sum += buffer_count * dt
        last_event_time = current_time

        if event_type == "arrival":
            total_arrivals += 1
            heapq.heappush(events,
                (current_time + gen_arrival(), "arrival", None))

            if busy_servers < n_servers:
                busy_servers += 1
                svc = gen_service()
                heapq.heappush(events,
                    (current_time + svc, "departure",
                     (current_time, current_time)))
            elif buffer_count < buffer_size:
                buffer_count += 1
                buffer_queue.append(current_time)
            else:
                total_rejected += 1

        elif event_type == "departure":
            total_served += 1
            if buffer_count > 0:
                buffer_count -= 1
                buf_arrival = buffer_queue.popleft()
                wait = current_time - buf_arrival
                total_wait_in_buffer += wait
                total_passed_through_buffer += 1
                svc = gen_service()
                heapq.heappush(events,
                    (current_time + svc, "departure",
                     (buf_arrival, current_time)))
            else:
                busy_servers -= 1

    total_sim_time = np.sum(state_time) or T
    P = state_time / total_sim_time

    Q     = total_served / total_arrivals if total_arrivals > 0 else 0.0
    A     = total_served / T
    P_rej = total_rejected / total_arrivals if total_arrivals > 0 else 0.0
    K     = busy_servers_sum / total_sim_time
    N_buf = buffer_sum / total_sim_time
    N_prog = sum(i * P[i] for i in range(n_states))
    T_prog = N_prog / A if A > 0 else 0.0
    T_buf  = (total_wait_in_buffer / total_passed_through_buffer
              if total_passed_through_buffer > 0 else 0.0)

    return {
        "P": P,
        "P0": P[0], "P1": P[1], "P2": P[2], "P3": P[3],
        "P4": P[4], "P5": P[5], "P6": P[6],
        "Q": Q, "A": A, "P_rej": P_rej, "K": K,
        "N_prog": N_prog, "T_prog": T_prog,
        "N_buf": N_buf, "T_buf": T_buf,
    }
\end{verbatim}

\end{document}
